\section{引言}
\subsection{研究背景与动机}
生成式人工智能（Gen-AI）正从一种小众的辅助工具演变为各行各业不可或缺的基础资源。随着其能力的快速增长，Gen-AI 不仅在优化工作流程，更在从根本上重塑职业的任务结构；相应地，高等教育也必须同步调整其培养体系，以回应“技能结构被技术重写”的现实挑战~\citep{mughal2025beyondautomationho,wong2025thefutureoflearnin}。

\subsection{问题重述}
首先题目的核心问题只有一个： post-secondary educational institutions如何更好的将 Gen-AI 融合进当前的培养方案。此核心问题分解为以下流程：\\首先，我们需要从 \textbf{STEM、Trade、Arts} 三个类别中各选一个具体职业，刻画这些职业在Gen-AI 影响下的未来变化。这里的“未来”可以被理解为要求确定一个明确的预测窗口，也可以视为一个广义的时间跨度。“当前轨迹”可解释为 Gen-AI 在该行业中的渗透率与采用率；“预期影响”则指这种渗透对岗位结构与从业者工作方式带来的可度量的变化。

其次，
在完成职业层面的分析后，赛题进一步要求：为每个职业各选定一个“对应的具体院校与具体项目”，并把建议落到这三个项目的具体改进措施上，要求用数学模型的输出为这些方案提供证指导性意见，而不是停留在原则性口号；这一点与近年的高教研究结论一致，即机构需要把 GenAI 的使用边界、学习收益与评估方式一起制度化，而非仅给出泛化口号~\citep{hamerman2024aninvestigationofg,jogezai2025generativeartifici}。   

最后，除了题目所给的建议，团队仍需自行思考额外的评价维度，并说明模型是否可以推广到其他机构或项目，需要明确写出可推广的条件与边界。